% =============================================================================
% ch05_fusion.tex — Chapter 5: Graph-SLAM: Factor Graphs, iSAM2 and g2o | NavigatorCrew
% Compiler: XeLaTeX | Language: Hebrew primary, English technical terms via \en{}
% =============================================================================

\selectlanguage{hebrew}

\section{גרף-\en{SLAM}: גרפי גורמים, \en{iSAM2} ו-\en{g2o}}
\label{sec:fusion}

% ── 5.1 Graph-SLAM Principles ───────────────────────────────────────────────
\subsection{עקרונות גרף-\en{SLAM} וניסוח מתמטי}
\label{subsec:graphslam_principles}

גרף-\en{SLAM} הפך לפרדיגמה הדומיננטית במערכות \en{SLAM} מודרניות,
והחליף את הגישות מבוססות \en{EKF} ברוב המערכות המסחריות והמחקריות.
התובנה המרכזית היא ניסוח \en{SLAM} כגרף גורמים דליל
(\en{Sparse Factor Graph}), שבו צמתים מייצגים תנוחות רובוט
ומיקומי נקודות ציון, וקשתות (גורמים) מייצגות אילוצים
מאודומטריה, תצפיות, סגירת לולאות ומידע קודם \cite{Thrun2005ProbRobotics}.

מטרת האופטימיזציה היא למצוא את התצורה בעלת ההסתברות המקסימלית
(\en{MAP}) של כל הצמתים בהינתן כל המדידות. בעיה זו מצטמצמת
לבעיית ריבועים פחותים לא-ליניארית (\en{Nonlinear Least Squares}),
הניתנת לפתרון באמצעות פירוק \en{Cholesky} דליל או פירוק \en{QR}.
שתי הספריות הדומיננטיות ליישום גרף-\en{SLAM} הן \en{g2o}
(\en{General Graph Optimization}) ו-\en{GTSAM}
(\en{Georgia Tech Smoothing and Mapping}).

\begin{equation}
    \mathbf{X}^* = \arg\min_{\mathbf{X}} \sum_{i} \| \mathbf{e}_{i}(\mathbf{X}) \|_{\Omega_i}^2
    \label{eq:graph_slam_objective}
\end{equation}

כאשר $\mathbf{X} = \{\mathbf{x}_1, \ldots, \mathbf{x}_T, \mathbf{l}_1, \ldots, \mathbf{l}_M\}$
הוא אוסף כל הצמתים (תנוחות ונקודות ציון), $\mathbf{e}_i$ היא פונקציית
השגיאה של גורם $i$, ו-$\Omega_i$ היא מטריצת האינפורמציה.

% ── 5.2 iSAM2: Incremental Updates ──────────────────────────────────────────
\subsection{\en{iSAM2}: עדכונים אינקרמנטליים באמצעות עץ בייס}
\label{subsec:isam2}

\en{iSAM2} (\en{Incremental Smoothing and Mapping using the Bayes tree})
מהווה פריצת דרך בתחום ההסקה האינקרמנטלית. במקום לפתור מחדש את
בעיית הריבועים הפחותים המלאה לאחר כל מדידה חדשה, \en{iSAM2}
מתחזק מבנה נתונים של עץ בייס (\en{Bayes Tree}) המאפשר עדכון
יעיל של רק החלק המושפע מהגרף. עץ הבייס הוא גרף מכוון חסר
מעגלים המקודד את התלויות המותנות של פירוק ה-\en{QR} של גרף
הגורמים.

האלגוריתם פועל בחמישה שלבים עיקריים: הוספת הגורם החדש לגרף,
זיהוי הקליקה (\en{Clique}) בעץ המושפעת מהמשתנים החדשים,
הסרת הקליקות המושפעות, ביצוע אלימינציה מחדש של המשתנים
המושפעים, והרכבה מחדש של העץ. תהליך זה מאפשר זמן אופטימיזציה
של 12 אלפיות שנייה על \en{KITTI 00}, לעומת 850 אלפיות שנייה
ל-\en{g2o} בקבוצות \cite{Thrun2005ProbRobotics}.

\begin{equation}
    p(\mathbf{X} | \mathbf{Z}) \propto \prod_{i} \phi_i(\mathbf{X}_i)
    \rightarrow \prod_{j} p(\mathbf{x}_j | \mathbf{S}_j)
    \label{eq:bayes_tree_conditional}
\end{equation}

משוואה~\ref{eq:bayes_tree_conditional} מציגה את המעבר מייצוג
גרף הגורמים (אגף שמאל) לייצוג עץ הבייס (אגף ימין), כאשר
$\mathbf{S}_j$ הם משתני האב של צומת $j$ בעץ.

% ── 5.3 Performance Comparison: iSAM2 vs g2o ────────────────────────────────
\subsection{השוואת ביצועים: \en{iSAM2} לעומת \en{g2o}}
\label{subsec:isam2_vs_g2o}

השוואת ביצועים על מערכי הנתונים \en{KITTI 00} ו-\en{TUM RGB-D}
מראה כי \en{iSAM2} ו-\en{g2o} משיגים דיוק דומה, אך \en{iSAM2}
מהיר משמעותית בעדכונים אינקרמנטליים. על \en{KITTI 00}, \en{iSAM2}
משיג \en{ATE RMSE} של 3.4 ס״מ לעומת 3.2 ס״מ ל-\en{g2o},
אך זמן האופטימיזציה הוא 12 אלפיות שנייה לעומת 850 אלפיות שנייה.
על \en{TUM RGB-D}, \en{iSAM2} משיג \en{ATE} של 1.9 ס״מ
בזמן של 4.5 אלפיות שנייה.

מבחינת צריכת זיכרון, \en{iSAM2} יעיל יותר: 198 \en{MB} לעומת
245 \en{MB} ל-\en{g2o} על \en{KITTI 00}. עם זאת, ל-\en{iSAM2}
יש מגבלות: הוא רגיש למדידות חריגות (\en{Outliers}), כגון סגירות
לולאה שגויות, העלולות להשחית את כל הגרף. טיפול בבעיה זו דורש
שימוש בפונקציות עלות \en{Robust} (כגון \en{M-Estimation}) או
\en{Dynamic Covariance Scaling} \cite{Grisetti2010g2o}.

\begin{table}[H]
\centering
\begin{tabular}{lccc}
\toprule
מערך נתונים & מדד & \en{iSAM2} & \en{g2o} \\
\midrule
\en{KITTI 00} & \en{ATE RMSE} [ס״מ] & 3.4 & 3.2 \\
\en{KITTI 00} & זמן אופטימיזציה [אלפיות ש׳] & 12 & 850 \\
\en{KITTI 00} & צריכת זיכרון [\en{MB}] & 198 & 245 \\
\en{TUM RGB-D} & \en{ATE RMSE} [ס״מ] & 1.9 & 2.1 \\
\en{TUM RGB-D} & זמן אופטימיזציה [אלפיות ש׳] & 4.5 & 320 \\
\bottomrule
\end{tabular}
\caption{השוואת ביצועים בין \en{iSAM2} ל-\en{g2o} על מערכי נתונים סטנדרטיים.}
\label{tab:graphslam_benchmark}
\end{table}

% ── 5.4 Robust Optimization ─────────────────────────────────────────────────
\subsection{אופטימיזציה עמידה בפני מדידות חריגות}
\label{subsec:robust_opt}

כדי להתמודד עם רגישות \en{iSAM2} למדידות חריגות, אנו משתמשים
בפונקציות עלות \en{Robust} המחליפות את פונקציית הריבועים
הרגילה. פונקציית ה-\en{Huber Loss} מוגדרת כ:

\begin{equation}
    \rho_{\delta}(e) =
    \begin{cases}
        \frac{1}{2} e^2 & |e| \leq \delta \\
        \delta(|e| - \frac{1}{2}\delta) & |e| > \delta
    \end{cases}
    \label{eq:huber_loss}
\end{equation}

כאשר $\delta$ הוא פרמטר הסף הקובע את המעבר משגיאה ריבועית
לשגיאה ליניארית. עבור $\delta = 1.0$, מדידות עם שגיאה גדולה
מ-1.0 מקבלות משקל מופחת, ובכך מוגבלת השפעתן על האופטימיזציה.

\begin{equation}
    \mathbf{X}^* = \arg\min_{\mathbf{X}} \sum_{i} \rho_{\delta}\left( \| \mathbf{e}_{i}(\mathbf{X}) \|_{\Omega_i} \right)
    \label{eq:robust_graph_slam}
\end{equation}

% ── 5.5 Summary ─────────────────────────────────────────────────────────────
\subsection{סיכום}
\label{subsec:graphslam_summary}

בפרק זה הצגנו את עקרונות גרף-\en{SLAM} ואת שתי הספריות
המובילות ליישומו: \en{iSAM2} ו-\en{g2o}. \en{iSAM2} מציע
ביצועים עדיפים בעדכונים אינקרמנטליים, אך דורש טיפול זהיר
במדידות חריגות באמצעות פונקציות עלות \en{Robust}. בפרק הבא
נתאר את אלגוריתם האודומטריה והמיפוי המבוסס \en{LiDAR}.